Before training the GP surrogate, we obtain a maximum \emph{a
posteriori} (MAP) estimate of the model parameters to serve as an
initial reference point for the MCMC samplers. The MAP estimate is
computed with \texttt{MaxLEV}, a maximum likelihood estimator for
\texttt{VPLanet} models that reads a JSON configuration file
specifying the free parameters, their bounds and priors, the
observational constraints, and the optimizer settings. For each
proposed parameter vector, \texttt{MaxLEV} writes \texttt{VPLanet}
input files, executes the forward model, parses the output, and
evaluates a Gaussian log-likelihood
\begin{equation}
  \ln\mathcal{L} = -\frac{1}{2}\sum_{k}
  \left(\frac{y_{k}^{\mathrm{obs}} - y_{k}^{\mathrm{model}}}
  {\sigma_{k}}\right)^{2},
  \label{eq:loglike}
\end{equation}
where $y_{k}^{\mathrm{obs}}$ and $\sigma_{k}$ are the observed value
and uncertainty of the $k$th constraint, and
$y_{k}^{\mathrm{model}}$ is the corresponding \texttt{VPLanet} output.
The two constraints used here are the bolometric luminosity
$L_{\mathrm{bol}}$ and the XUV-to-bolometric luminosity ratio
$\log(L_{\mathrm{XUV}}/L_{\mathrm{bol}})$.

The objective function---the negative log-posterior, incorporating
priors on stellar mass, $f_{\mathrm{sat}}$, age, and
$\beta_{\mathrm{XUV}}$---is minimized with differential evolution
\citep{StornPrice97} as implemented in \texttt{SciPy}
\citep{Virtanen20}. We use the \texttt{best1bin} strategy with a
population size of 3 per dimension, dithered mutation in $[0.5, 1.0]$,
a crossover probability of 0.7, and a convergence tolerance of 0.01,
running for up to 50 generations. The resulting MAP parameter vector
is used to initialize the \texttt{emcee} walkers and to verify that
the GP surrogate reproduces the true likelihood at the mode.
